\documentclass[]{article}
\usepackage[a4paper,includeheadfoot,margin=2.54cm]{geometry}

\usepackage[english,bulgarian]{babel}
\usepackage{polyglossia}
\usepackage{array}
\usepackage{amsmath}
\usepackage{amssymb}
\usepackage{graphicx}
\usepackage{float}
\usepackage{tabularx}
\usepackage[backend=biber,style=apa]{biblatex} % Use biblatex with biber
\usepackage{pdflscape}
\usepackage{adjustbox}
\usepackage{caption}
\usepackage{hyperref}

\hypersetup{
	colorlinks=true,      % Enable colored links
	linkcolor=black,     % Color of internal links
	citecolor=black,     % Color of citations
	filecolor=black,     % Color of file links
	urlcolor=black,      % Color of external links
	pdfborder={0 0 0},   % No visible border
	pdfinfo={
		Title={Predicting First‑Year Statistics Grades from Attendance, Honesty, and Quiz Results Cohort 2026/30},
		Author={R. Petrov},
		Subject={Decision Science - AEEM - FEBA - SU},
	}
}

%opening
\title{Predicting First‑Year Statistics Grades \\
	 from Attendance, Honesty, and Quiz Results\\
	 Cohort 2026/30}
\author{Riste Petrov}

\begin{document}
	
	\maketitle
	
	\begin{abstract}
		Sitting proximity to higher-achieving peers predicts improved exam performance among first-year students in an introductory statistics course. Observational and survey data from 49 students—attendance, quiz scores, seating location, and measured indicators of dishonesty—were analyzed using regression models to estimate peer effects on exams while controlling for individual covariates. Students seated near higher-performing classmates scored significantly higher on exams; this association persisted after adjusting for measured dishonesty and other covariates. Measured dishonesty correlated with lower performance, while attitudes toward mathematics were weak predictors. These results align with prior findings that local peer networks can influence classroom learning but extend them by explicitly measuring and ruling out detectable cheating as an alternative mechanism. The small sample limits causal claims and warrants replication; if robust, the findings suggest seating arrangements could be a low-cost intervention to boost early-course achievement.
	\end{abstract}
	
	\tableofcontents
	\newpage
	
	\section{Introduction}
	This research was motivated by R. Petrov’s appointment as Teaching Demonstrator for Statistics exercise sessions at Goce Delchev University — Shtip. During onboarding, formative questionnaires were introduced to assess incoming freshmen’s baseline knowledge and to encourage a more active, participatory learning environment. The process inspired development of a model to investigate possible causal relationships between students’ high-school background, academic honesty, attendance, homework completeness and tidiness, and their academic performance. The study combines questionnaire data and session observations to identify which background and behavioral factors most strongly predict success in introductory Statistics and to evaluate whether brief, targeted assessments can both reveal knowledge gaps and guide tailored instructional strategies.\\
	
	
	A substantial body of educational research links attendance, homework behaviors, academic integrity, and attitudes toward mathematics to student achievement, providing theoretical and empirical grounding for H0–H2. Regular class attendance consistently predicts higher achievement across grade levels—attendance improves exposure to instruction and opportunities for feedback, and correlates with course grades and test scores (Balfanz \& Byrnes, 2012; Gottfried, 2010). Homework completion and timeliness are robust predictors of learning: purposeful, well-designed homework improves mastery and exam performance, while missing or late assignments correlate with lower achievement and reflect weaker self-regulation (Cooper, Robinson, \& Patall, 2006; Trautwein et al., 2009). Studies that distinguish homework quantity from quality find that submission consistency and timely feedback matter more than mere volume, supporting inclusion of both numHW/additHW and lateHW in the models (König et al., 2020).\\
	
	Academic dishonesty complicates interpretation of measured performance: dishonest behaviors can inflate short-term scores for some students while undermining genuine learning and classroom norms; indices of dishonesty relate negatively to subsequent learning outcomes when long-term mastery is the target (McCabe, Butterfield, \& Treviño, 2012; Anderman \& Murdock, 2007). This literature motivates modeling DISS both as a main predictor and in interactions with seating proximity variables to detect whether localized peer effects reflect collaboration versus collusion. Attitudes toward mathematics (affect and self-efficacy) exert consistent but often smaller direct effects on achievement than proximal behavioral factors; positive attitudes predict persistence and strategy use and typically operate indirectly through engagement and homework behaviors (Pajares \& Graham, 1999; Zimmerman, 2000), justifying testing whether attitMATH adds explanatory power beyond attendance and homework in H0.\\
	
	Peer effects and seating proximity have empirical support in education economics and social psychology: proximity to higher-achieving peers can raise individual achievement via social learning and modeling, while clustering with lower-achieving peers may depress outcomes; experimental and quasi-experimental classroom studies show measurable peer spillovers but with context-dependent effect sizes, which motivates separate PROXI25 and PROXI75 indicators and interaction tests with DISS (Sacerdote, 2011; Carrell, Sacerdote, \& West, 2013). Finally, early low‑stakes assessments (diagnostic quizzes) reliably forecast later exam performance and, when followed by targeted remediation, improve outcomes—supporting inclusion of firstQIZZ and autoregressive links between Y1 and Y2 (Black \& Wiliam, 1998; Hattie \& Timperley, 2007). Overall, extant evidence indicates that proximal behavioral predictors (attendance, homework habits, early quizzes, and academic honesty) often explain variance in achievement beyond attitudinal measures, while peer proximity and interactions can moderate these relationships—justifying the functional forms and interaction terms specified for H0–H2.\\
	
	
	\subsection{Hypotheses}
	
	\subsubsection{H0 - Attendance, homework completion/timeliness, and academic
		honesty predict test performance more strongly than attitude toward
		math}
		The sample was formed out of 49 regular (non-test only) students which have been unubstorcatble have been observed and their beahviour was noted
		The functional forms for the hypothesis is as follows for both midterms ($Y_1 \& Y_2$) :
		
		$$
		\begin{aligned}
			Y_1 &= \alpha + \beta_1 \cdot \text{GENDER} + \beta_2 \cdot Y_1\text{GROUP} + \beta_3 \cdot \text{averATND}_1 + \beta_4 \cdot \text{numHW}_1 \\
			&\quad + \beta_5 \cdot \text{paperHW}_1 + \beta_6 \cdot \text{lateHW}_1 + \beta_7 \cdot \text{additHW}_1 + \beta_8 \cdot \text{DISS}_1 + \beta_9 \cdot \text{atitMATH}
		\end{aligned}
		$$
		$$
		\begin{aligned}
			Y_2 &= \alpha + \beta_1 \cdot \text{GENDER} + \beta_2 \cdot Y_2\text{GROUP}  + \beta_3 \cdot \text{averATND}_2 +   \\
			&\quad + \beta_4 \cdot \text{numHW}_2 + \beta_5 \cdot \text{additHW}_2 + \beta_6 \cdot \text{DISS}_2 + \beta_7 \cdot \text{atitMATH}
		\end{aligned}
		$$
		The functional form for the second midterm is reduced due to the fact that the students had one homework and one additional homework assignment and the potential lack of variance could be  cause of issues such as multicolinearity.
		
	
	\subsubsection{H1 - Seating proximity to lower- or higher-achieving peers effects student performance compared with students near the median.}
		
		The baseline functional forms of the hypothesis are as follows for both midterms ($Y_1 \& Y_2$) :
		$$
		\begin{aligned}
			Y_1 &= \alpha + \beta_1 \cdot \text{GENDER} + \beta_2 \cdot Y_1\text{GROUP} + \beta_3 \cdot \text{PROXI25}_1 + \beta_4 \cdot \text{PROXI75}_1 \\
		\end{aligned}
		$$
		$$
		\begin{aligned}
			Y_2 &= \alpha + \beta_1 \cdot \text{GENDER} + \beta_2 \cdot Y_2\text{GROUP} + \beta_3 \cdot \text{PROXI25}_2 + \beta_4 \cdot \text{PROXI75}_2 \\
		\end{aligned}
		$$
		
		
		In order to single out cheating-endeavors and to test a similar hypothesis to the previous one a sub-hypothesis was formed to test the effects of the dishonesty (DISS) index and the joint effects with the sitting proximity to lower and higher achieving students: 
		
		$$
		\begin{aligned}
			Y_1 &= \alpha + \beta_1 \cdot \text{GENDER} + \beta_2 \cdot Y_1\text{GROUP} + \beta_3 \cdot \text{DISS}_1    \\
			&\quad + \beta_4 \cdot \text{PROXI25}_1 + \beta_5 \cdot \text{PROXI75}_1  + \beta_6 \cdot \text{INT}(\text{DISS}_1,\text{PROXI75}_1)
		\end{aligned}
		$$
		
		$$
		\begin{aligned}
			Y_2 &= \alpha + \beta_1 \cdot \text{GENDER} + \beta_2 \cdot Y_2\text{GROUP} + \beta_3 \cdot \text{DISS}_2    \\
			&\quad + \beta_4 \cdot \text{PROXI25}_2 + \beta_5 \cdot \text{PROXI75}_2  + \beta_6 \cdot \text{INT}(\text{DISS}_2,\text{PROXI75}_2)
		\end{aligned}
		$$
		
		The interaction between PROXI25 and DISS was omitted due to lack of a possible causal link of malicious intent of sitting next to a lower achieving student with the possibility to cheat from them. \\
		
	\subsubsection{H2 - Lower attendance, late homework submission, digitally completed homework, low quiz scores, and relatively hard perceived homework are associated with lower final results, 	and the opposite patterns predict higher results.}
		
		The functional form for the both midterms is as follows ($Y_1$ \& $Y_2$):
		$$
		\begin{aligned}
			Y_1 &= \alpha + \beta_1 \cdot \text{GENDER} + \beta_2 \cdot Y_1\text{GROUP} + \beta_3 \cdot \text{averATND\_H}_1   \\
			&\quad + \beta_4 \cdot \text{lateHW}_1 + \beta_5 \cdot \text{paperHW}_1  + \beta_6 \cdot \text{averQIZZ\_H}_1 + \beta_7 \cdot \text{hardHW\_H}_1
		\end{aligned}
		$$
		
		$$
		\begin{aligned}
			Y_2 &= \alpha + \beta_1 \cdot \text{GENDER} + \beta_2 \cdot Y_2\text{GROUP} + \beta_3 \cdot \text{averATND\_H}_2   \\
			&\quad + \beta_4 \cdot \text{lateHW}_2 + \beta_5 \cdot \text{paperHW}_2  + \beta_6 \cdot \text{averQIZZ\_H}_2 + \beta_7 \cdot \text{hardHW\_H}_2
		\end{aligned}
		$$
	
	\subsubsection{Notes for hypotheses}
	
	Additionally, to test autoregressive effects, the variables firstQIZZ and $Y_1$ were included in the models for the two midterms in order to assess whether the previous test exhibits a lagged autoregressive effect. Also the use of two-stage least squares (2SLS) was used, treating the prior exam score as endogenous and instrumenting it appropriately.
	
	\subsection{Sample}
	The total sample comprised 49 regular (non–exam‑only) students enrolled in the mandatory course Statistics for Economists (4EF113421) across all five majors in the Faculty of Economics. Only 30 of these students completed the first exercise‑session knowledge assessment (quiz). \\
	
	This small sample size limits the reliability and robustness of regression estimates. The experimental design and data granularity prevented pooling additional student groups to validate or refute the proposed hypotheses, which increases the risk of low statistical power and unstable parameter estimates.\\
	
	The list of baseline variables collected:
	\begin{itemize}
		\item [] name of variable - type - description
		\item e26.02.23(date) - logical - evidence list attendance.
		\item eindex26.02.23 (date) - continuous - index of evidence list
		\item honesty-evidence list - logical - one lecture after the actual lecture, the list is checked to see if the absent person signs being checked when he/she was absent.
		\item honesty-homework - logical - every lecture there is a quiz which asks the participants if they have written homework.
		\item q26.02.23(date) - logical - quiz taken or not taken.
		\item \%26.02.23(date) - continuous [0,1] - the \% score achieved at the quiz for that lecture.
		\item 26.02.23(date)index - continuous - the index of the paper list of attendance made each lecture.
		\item p26.02.23(date) - logical - present or absent on the paper attendance list.
		\item h26.02.23(date) - logical - the student has written the homework assigned for the date lecture.
		\item paper - logical - the submitted homework has been written by hand (not online).
		\item late - logical - the submitted homework has been submitted late after the deadline.
		\item Homework-hardness - categorical [1,6] - 1 is equal to impossible and 6 is equal to easy.
	\end{itemize}
	\vspace{1em}
	
	To address the high dimensionality and irregular structure of the data, the data was aggregated related baseline variables into summary measures (means and indices) where appropriate to simplify interpretation and reduce the number of predictors. This aggregation was motivated by the hybrid time–cross-sectional design with two non-equivalent midterm periods and inconsistent measurement timing, which precluded constructing a balanced panel or pooling additional cohorts. Consequently, the analyses use these aggregated summary measures and concentrate on descriptive robustness checks while presenting evidence suggestive of causal relationships; nevertheless, the small effective sample and irregular data structure still limit statistical power and the stability of multivariable regression estimates.
	
	\newpage
	
	\subsection{Experimental Design}
	
	To preserve natural student behavior, experimental procedures were not disclosed. Each quiz was framed as an opportunity for students to do their best and participants were explicitly told that quiz results would not affect final grades. Quizzes consisted of multiple‑choice items drawn directly from lecture material and exercises to reflect taught content. Quizzes were administered using Google Forms because funds for paid platforms (e.g., Kahoot) were unavailable; those platforms would have allowed timed responses and gamification. The first quiz collected demographic‑style information; the final dataset is anonymous. (using annonINDX for privacy reasons) The demographic questionnaire translated into English is shown below. \\
	\begin{enumerate}
		\item Number of index - format (name.number of index) Example: example.00000
		\item I graduated:
		\subitem General High School
		\subitem Economics / Law Specialized High School
		\subitem Other non-related Specialized High School
		\item My major is:
		\subitem Finance
		\subitem Business Economics
		\subitem Marketing \& Management
		\subitem Accounting \& Audit
		\subitem International economics \& business
		\subitem Other
		\item I consider Mathematics as:
		\subitem scale awful 1 - 6 favorite
		\item Question about passing grade of Mathematics was added but the results were not available yet at the time of doing the quiz. (possible reporting collaboration with the Mathematics professor could ensure better results, possible updated version of the research)
		\subitem scale F 0 - 5 A.
	\end{enumerate}

	
	The same non‑disclosure applied to the two attendance lists, with the addition of a quiz‑participant list used to verify actual presence; requiring consistency across three independent records made coordinated manipulation substantially harder. Below there is a presentation of the intro quiz questions non material related:
	
	\begin{enumerate}
		\item Number of index - format (name.number of index) Example: example.00000
		\item Did i write homework?
		\subitem Yes or No.
		\item How hard was the homework? 
		\subitem scale Impossible 1 - 6 To easy
	\end{enumerate}
	The defaults of the questions were Yes and 3 in order to have a non-responsive trap answer, and to gauge if the respondents payed attention or were honest about the answers.
	
	\subsection{Feature Engineering}
	As touched upon about in the subsection about the Sample the "raw" collected data was averaged / indexed in order to derived more usable variables.\\
	
	The full list of used variables is:
	
	\begin{itemize}
		\item [] name of variable - midterm 1. (type) / midterm 2. (type) - description
		\subitem types = [cont(ious), categ(orical),log(ical)]
		\item anonINDX - cont - index starting from 1 formed by using the index of the student index number in a ascending order.
		\item GENDER - log - male (0) / female (1) assumed based on the name of each student. \begin{small}(not questioned at all)\end{small}
		\item HSCH - categ [1-3] - type of high-school, self-reported at the firstQIZZ. 
		\subitem types = [1 - general high school, 2 - economics/law specialized , 3 - other non-related specialized]
		\item MJR - categ [1-6] - type of major enrolled, self-reported at the firstQIZZ.
		\subitem types = [1 - Finance, 2 - Business Economics , 3 - Marketing \& Management, 4 - Accounting \& Audit, 5 - International economics \& business, 6 - Other]
		\item atitMATH - categ [1-6] - attitude towards Mathematics, self-reported at the firstQIZZ.
		\subitem scale = [1=awful - 6=favorite]
		
		\item averQIZZ - 1 (cont) / 2 (cont) - scaled [0,1] average in \% points for uniformity from the quizzes taken,
		\item diffQIZZ - 1 (cont) / 2 (cont) - scaled [0,1] difference of between the first and the average in \% points for uniformity from the quizzes taken, 
		\item averATND - 1 (cont) / 2 (cont) - average attendance calculated using: 
		\subitem numP - 1 (cont) / 2 (cont) - total number paper list attendance.
		\subitem numE - 1 (cont) / 2 (cont) - total number evidence list attendance.
		\subitem numQ - 1 (cont) / 2 (cont) - total number of quizzes taken.
		\item DISS - 1 (cont) / 2 (cont) - dishonesty index compiled using:
		\subitem honLST - log - NOR statement logic (once dishonest, forever dishonest), resulting in 3 points for falsifying being on a lecture with a signature on which the student was absent.
		\subitem honHW - log - NOR statement logic (once dishonest, forever dishonest), resulting in 1 point for clumsy answer or lying about writing homework.
		\subitem $Y_1$ CHEAT \& $Y_2$ CHEAT - log - suspected cheating during the midterm exam, resulting in 2 points for a behavior that indicates a possible cheating attempt.\\
		
		\vspace{0.5em}
		Notes:
		 \subitem The amount of DISS after Midterm 1 (2DISS) equals the initial DISS score up to Midterm 1 (1DISS) plus any additional changes observed thereafter — i.e., it records the combination of "old" and "new" behavior.\\
		 \subitem CHEAT denotes observed behaviors during a test that may indicate possible cheating but are not sufficiently serious or conclusive to trigger the university's standard immediate sanctions (automatic fail and one‑year ban on retaking the midterm/exam).\\
		
		\item Homework related variables:
		\subitem numHW - 1 (cont) / 2 (log) - the number of homework written during the parts of the semester.
		\subitem paperHW - 1 (cont) / 2 (log) - the average rate of submitting homework written by hand or electronic.
 		\subitem lateHW - 1 (cont) / 2 (log) - the average rate of submitting homework late.
		\subitem hardHW - 1 (categ) / 2 (categ) - the average of rate of hardness of the homework, scaled [0,1] from the original scale [1-6].
		\item Additional Homework (Course Project) related variables:
		\subitem additHW - 1 (log) / 2 (log) - measuring if additional homework was submitted or not. 
		\subitem lateADHW - 1 (log) - late submission was available only for the first midterm additional homework.
		\subitem electADHW - 2 (log) - electronic submission were measured only for the second midterm additional homework. (as a substitute for the late submission)
		\item Sitting proximity variables
		\subitem PROXI25 - 1 (cont) / 2 (cont) - average rate of sitting near (maximum distance 2 index spaces) next to a student having a score smaller than the 25th Percentile
		\subitem PROXI75 - 1 (cont) / 2 (cont) - average rate of sitting near (maximum distance 2 index spaces) next to a student having a score bigger than the 75th Percentile
		\subitem Note:
		\subsubitem The distance and sitting was measured by the (date)index - raw variable which is reconstructed at the start of each exercise session when students passed the list to submit attendance. This index captures clustering by reflecting students’ seating patterns. For proximity calculations, the two index positions immediately ahead and behind each student were used. The classroom consists of rows of two benches (four seats per connected unit), so this window provides a reasonable approximation of typical neighboring seats. To reduce the influence of occasional anomalies, each student’s average seating proximity weighs out outlier exercise sessions.\\
		
		Depended variable
		\vspace{1em}
		\item $Y_1$ \& $Y_2$ - cont [0,1] - in \% points, both 0-20.
		 \subitem $Y_1$GROUP \& $Y_2$GROUP are also available as controls such as GENDER variable.  
		 \subitem Notes:
		 \subsubitem $Y_1$ was made easier relative to the taught material and there was a possibility of cheating. Assessment composition: 5 points multiple‑choice, 5 points fill‑in, 10 points from an averaging exercise.
		 \subsubitem $Y_2$ was one grade more difficult relative to the taught material and cheating was effectively prevented by dynamic assignment elements. Assessment composition: 15 points from an OLS dynamic exercise, 5 points from a dynamic productivity‑ratios and indexing exercise.
		 \subsubitem hard$Y_1$ was measured after the midterm during the first subsequent exercise session.
		 
 	\end{itemize}
	
	
	All variables continuous variables are available in a centered form which was made as a remedy in order to deal with multicolinearity and similar econometric issues. The use of variable\_H is a logical(binary) variable in which the 1 is bigger than the average of the sample.  
	
	
	\section{Results}
	
	\subsection{Modeling firstQIZZ based on demographic-like data}
	A potential model of the demographic-like data was formed in the following functional form:
	
	$$
	\begin{aligned}
		\text{firstQIZZ} &= \alpha + \beta_1 \cdot \text{GENDER} + \beta_2 \cdot \text{anonINDX}   \\
		&\quad + \beta_3 \cdot \text{HSCH}  + \beta_4 \cdot \text{MJR} + \beta_5 \cdot \text{ATITMATH}
	\end{aligned}
	$$
	Below the output of the OLS model:
	\begin{center}
		\begin{tabular}{c|cccc}
			\hline
			Depended Variable : firstQIZZ & where & N = 25 & out of & 49 \\
			\hline
			Variable & Coefficient & Std. Error & t-Statistic & Prob.   \\
			\hline
			C & 0.259409 & 0.121589 & 2.133481 & 0.0461 \\
			ANONINDX & 0.00287 & 0.002443 & 1.174646 & 0.2546 \\
			GENDER & 0.098091 & 0.070665 & 1.388116 & 0.1812 \\
			HSCH & 0.005304 & 0.042282 & 0.125454 & 0.9015 \\
			MJR & -0.042932 & 0.031729 & -1.353102 & 0.1919 \\
			ATITMATH & -0.005642 & 0.022588 & -0.249773 & 0.8054 \\
			\hline
			R-squared & 0.202733 &  &  &  \\
			Adjusted R-squared & -0.007074 &  &  &  \\
			Akaike info criterion & -0.779252 &  &  &  \\
			Schwarz criterion & -0.486722 &  &  &  \\
			Log likelihood & 15.74065 &  &  &  \\
			F-statistic & 0.966282 &  &  &  \\
			Prob(F-statistic) & 0.462877 &  &  &  \\
			\hline
		\end{tabular}
	\end{center}
	The F‑statistic and p‑value indicate the model is not statistically significant; this may reflect cheating and/or omitted variables — most notably prior mathematics achievement (whether the student passed Math for Economists course and the grade received) — which could be causally related to firstQIZZ performance and influence the model’s estimates.
	
	\subsection{H0 - Honesty vs Attitude towards Math}
	The regression results indicate a robust, substantive negative association between the dishonesty index (DISS) and student test performance: a one‑point increase in DISS is associated with an average decline of approximately 7.6 to 11.2 percentage points in test scores across models. By contrast, the attitMATH measure is consistently statistically insignificant and does not predict performance.\\
	
	Some other covariates show occasional significance (numHW, averATND), but these effects are neither consistent nor as persistent as the DISS effect; averATND even appears with the opposite sign to that hypothesized. Control variables such as GENDER and Y$_{1or2}$GROUP are sometimes significant but should not be interpreted causally given sample imbalance across gender and group. \\
	
	The dishonesty index (DISS) is constructed by summing three components: honLST (signature honesty on the evidence list; 3 points), Y$_{1or2}$CHEAT (cheating indicator; 2 points), and honHW (self‑reported homework honesty; 1 point).
	
	 \begin{landscape}
	 	\begin{table}[htp] 
	 		\centering 
	 		\begin{adjustbox}{height=0.525\textheight,width=1.5\textwidth}
	 		\begin{tabular}{c|ccc|ccc|ccc}
		\hline
		Model Name & 1H0 & 1H0\_ar & 1H0\_ar\_2sls & 1H0\_reduced & 1H0\_reduced\_ar & 1H0\_reduced\_ar\_2sls & 2H0 & 2H0\_ar & 2H0\_ar\_2sls \\
		\hline
		Depended Variable : & $Y_1$ & $Y_1$ & $Y_1$ & $Y_1$ & $Y_1$ & $Y_1$ & $Y_2$ & $Y_2$ & $Y_2$ \\
		Sample Size (N)& 26 & 24 & 24 & 27 & 25 & 25 & 27 & 27 & 24  \\
		\hline
		Endogenous &  &  & $\hat{firstQIZZ}$ &  &  & $\hat{firstQIZZ}$ &  &  & $\hat{Y_1}$ \\
		\hline
		C & 0.9994 & 0.9175 & 0.9470 & 0.9123 & 0.8429 & 0.7811 & 0.1203 & -0.1095 & -0.0942 \\
		& 0.1908 & 0.1982 & 0.2201 & 0.1885 & 0.1901 & 0.2494 & 0.3583 & 0.3533 & 0.7545 \\
		& 0.0001*** & 0.0005*** & 0.0009*** & 0.0001*** & 0.0004*** & 0.0064 & 0.7407 & 0.7602 & 0.9023 \\
		\hline
		firstQIZZ or $Y_1$ &  & 0.0215 & 0.3689 &  & -0.1839 & 0.3892 &  & 0.6230 & 1.4642 \\
		&  & 0.2256 & 0.5833 &  & 0.1847 & 0.6534 &  & 0.3150 & 0.8655 \\
		&  & 0.9255 & 0.5381 &  & 0.3343 & 0.5597 &  & 0.0635* & 0.1114 \\
		\hline
		GENDER & 0.1162 & 0.1153 & 0.1137 & 0.0973 & 0.0971 & 0.0794 & 0.1440 & 0.0113 & -0.2313 \\
		& 0.0582 & 0.0569 & 0.0619 & 0.0575 & 0.0569 & 0.0744 & 0.0840 & 0.1031 & 0.2455 \\
		& 0.0631* & 0.0636* & 0.089* & 0.1069 & 0.1069 & 0.3020 & 0.1026 & 0.9143 & 0.3611 \\
		$Y_{1or2}$GROUP & -0.1349 & -0.1356 & -0.1182 & -0.1466 & -0.1443 & -0.1482 & 0.0563 & 0.0201 & -0.1581 \\
		& 0.0537 & 0.0603 & 0.0707 & 0.0527 & 0.0527 & 0.0668 & 0.0764 & 0.0735 & 0.3141 \\
		& 0.0232** & 0.0423** & 0.1183 & 0.0118** & 0.0146** & 0.0414** & 0.4701 & 0.7876 & 0.6221 \\
		\hline
		averATND$_{1or2}$ & -0.1212 & -0.0927 & -0.1313 & -0.0896 & -0.0497 & -0.0730 & -0.0872 & -0.0146 & 0.0115 \\
		& 0.0469 & 0.0606 & 0.0884 & 0.0453 & 0.0517 & 0.0700 & 0.0748 & 0.0787 & 0.1662 \\
		& 0.02** & 0.1504 & 0.1613 & 0.0623* & 0.3507 & 0.3127 & 0.2581 & 0.8550 & 0.9457 \\
		numHW$_{1or2}$ & -0.0265 & -0.0213 & -0.0576 & 0.0025 & -0.0057 & 0.0047 & 0.2474 & 0.2234 & 0.3401 \\
		& 0.0584 & 0.0772 & 0.1006 & 0.0384 & 0.0405 & 0.0525 & 0.0937 & 0.0881 & 0.2317 \\
		& 0.6563 & 0.7871 & 0.5766 & 0.9483 & 0.8903 & 0.9293 & 0.0161** & 0.0207** & 0.1628 \\
		paperHW$_{1or2}$ & 0.1073 & 0.0773 & 0.1829 &  &  &  &  &  &  \\
		& 0.1155 & 0.1482 & 0.2276 &  &  &  &  &  &  \\
		& 0.3665 & 0.6106 & 0.4362 &  &  &  &  &  &  \\
		lateHW$_{1or2}$ & 0.2414 & 0.2749 & 0.3808 &  &  &  &  &  &  \\
		& 0.1406 & 0.1522 & 0.2311 &  &  &  &  &  &  \\
		& 0.1053 & 0.0941* & 0.1234 &  &  &  &  &  &  \\
		additHW$_{1or2}$ & 0.0134 & 0.0054 & -0.0817 & 0.0413 & 0.0833 & -0.0432 & 0.3869 & 0.2743 & 0.1499 \\
		& 0.0922 & 0.1056 & 0.1755 & 0.0919 & 0.0968 & 0.1821 & 0.2049 & 0.1991 & 0.5534 \\
		& 0.8864 & 0.9601 & 0.6492 & 0.6582 & 0.4023 & 0.8155 & 0.0743* & 0.1852 & 0.7902 \\
		DISS$_{1or2}$ & -0.0910 & -0.0951 & -0.0834 & -0.0767 & -0.0891 & -0.0621 & -0.1120 & -0.0825 & -0.0320 \\
		& 0.0272 & 0.0274 & 0.0348 & 0.0266 & 0.0268 & 0.0445 & 0.0408 & 0.0409 & 0.0756 \\
		& 0.0041*** & 0.0042*** & 0.0322** & 0.0096*** & 0.0043*** & 0.1815 & 0.013** & 0.0587* & 0.6785 \\
		\hline
		atitMATH & 0.0199 & 0.0160 & 0.0306 & 0.0202 & 0.0115 & 0.0271 & 0.0244 & 0.0078 & -0.0246 \\
		& 0.0174 & 0.0238 & 0.0341 & 0.0175 & 0.0204 & 0.0306 & 0.0275 & 0.0269 & 0.0535 \\
		& 0.2704 & 0.5126 & 0.3860 & 0.2623 & 0.5790 & 0.3894 & 0.3858 & 0.7755 & 0.6529 \\
		\hline
		$R^2$ & 0.7098 & 0.7397 & 0.6922 & 0.6585 & 0.6897 & 0.5030 & 0.5740 & 0.6500 & 0.3038 \\
		Adjusted $R^2$ & 0.5466 & 0.5395 & 0.4554 & 0.5327 & 0.5346 & 0.2545 & 0.4170 & 0.4945 & -0.0675 \\
		Akaike info criterion & -1.2957 & -1.3597 & Not calculated & -1.2834 & -1.3566 & Not calculated & -0.2715 & -0.3941 & Not calculated \\
		Schwarz criterion & -0.8118 & -0.8198 & Not calculated & -0.8995 & -0.9178 & Not calculated & 0.1125 & 0.0379 & Not calculated \\
		Log likelihood & 26.8439 & 27.3168 & Not calculated & 25.3263 & 25.9577 & Not calculated & 11.6652 & 14.3200 & Not calculated \\
		\hline
		F-statistic & 4.3481 & 3.6941 & 3.1633 & 5.2337 & 4.4457 & 2.7426 & 3.6568 & 4.1790 & 1.9622 \\
		Prob(F-statistic) & 0.005202*** & 0.01527** & 0.027619** & 0.00187*** & 0.005414*** & 0.040884** & 0.011441** & 0.005607*** & 0.1240 \\
		J-statistic & NA & NA & 2.8038 & NA & NA & 0.9201 & NA & NA & 0.4597 \\
		Prob(J-statistic) & NA & NA & 0.2461 & NA & NA & 0.6313 & NA & NA & 0.7947 \\
		\hline
	\end{tabular}
	\end{adjustbox}
	\caption{Regression analysis for H0}  
	\end{table} 
	\end{landscape}
	
	\subsection{H1 - Sitting proximity \& Honesty}
	Proximity to higher‑achieving peers (PROXI75) has a robust, substantively meaningful positive effect on end‑of‑course scores across multiple specifications. Estimates are statistically significant in OLS and when controlling for dishonesty (DISS), and the result remains in some of the autoregressive and 2SLS models, indicating the effect is not driven by measured dishonesty or by endogeneity in peer assignment.\\
	
	In contrast, proximity to lower‑achieving peers (PROXI25) yields no consistent or significant effect in any specification, providing no support for the idea that sitting near lower‑achieving classmates materially alters performance.\\
	
	Two findings reduce concerns that the PROXI75 effect reflects cheating: (1) the coefficient on PROXI75 remains after controlling for DISS, so students with lower honesty scores do not explain the result; and (2) similar PROXI75 estimates appear for the second midterm, when opportunities for collusion were much smaller — consistent with mechanisms such as knowledge transfer, social learning, or increased effort rather than copying.\\
	
	Overall, the results support the hypothesis that sitting next to higher‑achieving peers improves outcomes, while offering no evidence that proximity to lower‑achieving peers affects performance. Tables 2 and 3 present the model estimates.
	
	\begin{landscape}
		\begin{table}[htp] 
			\centering 
			\begin{adjustbox}{height=0.55\textheight,width=1.6\textwidth}
				\begin{tabular}{c|ccc|ccc|ccc}
					\hline
					Model Name & 1h1 & 1h1\_ar & 1h1\_ar\_2sls & 1h1\_diss & 1h1\_diss\_ar & 1h1\_diss\_ar\_2sls & 1h1\_diss\_int & 1h1\_diss\_int\_ar & 1h1\_diss\_int\_ar\_2sls \\
					\hline
					Sample Size (N) & 46 & 25 & 25 & 46 & 25 & 25 & 46 & 25 & 25 \\
					\hline
					Endogenous &  &  & $\hat{firstQIZZ}$ &  &  & $\hat{firstQIZZ}$ &  &  & $\hat{firstQIZZ}$ \\
					\hline
					C & 0.4553 & 0.5070 & 0.5102 & 0.4809 & 0.7030 & 0.6605 & 0.4182 & 0.7466 & 0.7193 \\
					& 0.1209 & 0.1839 & 0.1865 & 0.1303 & 0.1667 & 0.1864 & 0.1575 & 0.1821 & 0.2056 \\
					& 0.0005*** & 0.0126** & 0.0131** & 0.0007*** & 0.0005*** & 0.0023*** & 0.0114** & 0.0007*** & 0.0028*** \\
					\hline
					firstQIZZ &  & -0.0406 & -0.0873 &  & -0.1309 & 0.1924 &  & -0.1142 & 0.2555 \\
					&  & 0.2024 & 0.4746 &  & 0.1717 & 0.3476 &  & 0.1764 & 0.3718 \\
					&  & 0.8431 & 0.8560 &  & 0.4557 & 0.5867 &  & 0.5260 & 0.5013 \\
					\hline
					GENDER & 0.0950 & 0.1176 & 0.1197 & 0.0922 & 0.1082 & 0.0948 & 0.0981 & 0.1085 & 0.0937 \\
					& 0.0621 & 0.0870 & 0.0893 & 0.0628 & 0.0727 & 0.0805 & 0.0637 & 0.0739 & 0.0839 \\
					& 0.1334 & 0.1923 & 0.1958 & 0.1497 & 0.1542 & 0.2541 & 0.1316 & 0.1604 & 0.2794 \\
					\hline
					Y$_1$GROUP & -0.1079 & -0.1224 & -0.1179 & -0.1103 & -0.1400 & -0.1684 & -0.1000 & -0.1436 & -0.1771 \\
					& 0.0408 & 0.0614 & 0.0742 & 0.0414 & 0.0517 & 0.0621 & 0.0441 & 0.0528 & 0.0657 \\
					& 0.0116** & 0.0608* & 0.1284 & 0.011** & 0.0143** & 0.0143** & 0.0288** & 0.0145** & 0.0154** \\
					\hline
					DISS &  &  &  & -0.0115 & -0.0770 & -0.0687 & 0.0116 & -0.1107 & -0.1172 \\
					&  &  &  & 0.0206 & 0.0254 & 0.0287 & 0.0382 & 0.0578 & 0.0651 \\
					&  &  &  & 0.5804 & 0.0071*** & 0.0279** & 0.7634 & 0.0725* & 0.0895* \\
					PROXI$_{25}$ & 0.0039 & 0.0125 & 0.0087 & 0.0050 & -0.0073 & 0.0203 & 0.0054 & 0.0022 & 0.0375 \\
					& 0.0948 & 0.1324 & 0.1371 & 0.0957 & 0.1108 & 0.1238 & 0.0962 & 0.1136 & 0.1309 \\
					& 0.9674 & 0.9261 & 0.9503 & 0.9584 & 0.9480 & 0.8716 & 0.9557 & 0.9849 & 0.7777 \\
					PROXI$_{75}$ & 0.1421 & 0.0597 & 0.0650 & 0.1404 & 0.0316 & -0.0008 & 0.2114 & -0.0418 & -0.1124 \\
					& 0.0650 & 0.1009 & 0.1120 & 0.0656 & 0.0848 & 0.0973 & 0.1188 & 0.1418 & 0.1701 \\
					& 0.0345** & 0.5606 & 0.5685 & 0.0386** & 0.7141 & 0.9933 & 0.083* & 0.7717 & 0.5175 \\
					INT(DISS, PROXI$_{75}$) &  &  &  &  &  &  & -0.0386 & 0.0529 & 0.0777 \\
					&  &  &  &  &  &  & 0.0537 & 0.0812 & 0.0935 \\
					&  &  &  &  &  &  & 0.4763 & 0.5234 & 0.4177 \\
					\hline
					R\^{}2 & 0.3578 & 0.4466 & 0.4450 & 0.3628 & 0.6340 & 0.5619 & 0.3711 & 0.6429 & 0.5506 \\
					Adjusted R\^{}2 & 0.2952 & 0.3009 & 0.2990 & 0.2831 & 0.5120 & 0.4158 & 0.2744 & 0.4958 & 0.3656 \\
					Akaike info criterion & -1.0326 & -1.0180 & Not calculated & -0.9968 & -1.3514 & Not calculated & -0.9665 & -1.2960 & Not calculated \\
					Schwarz criterion & -0.8338 & -0.7255 & Not calculated & -0.7583 & -1.0101 & Not calculated & -0.6882 & -0.9060 & Not calculated \\
					Log likelihood & 28.7490 & 18.7248 & Not calculated & 28.9269 & 23.8923 & Not calculated & 29.2299 & 24.2005 & Not calculated \\
					\hline
					F-statistic & 5.7119 & 3.0664 & 3.0566 & 4.5548 & 5.1960 & 4.3112 & 3.8361 & 4.3720 & 3.4944 \\
					Prob(F-statistic) & 0.00095*** & 0.033978** & 0.03437** & 0.002216*** & 0.002924*** & 0.007239*** & 0.004221*** & 0.006108*** & 0.016547** \\
					J-statistic & NA & NA & 4.9859 & NA & NA & 3.9408 & NA & NA & 3.0710 \\
					Prob(J-statistic) & NA & NA & 0.1728 & NA & NA & 0.2679 & NA & NA & 0.3808 \\
					\hline
				\end{tabular}
			\end{adjustbox}
			\caption{Regression analysis for H1 for the first midterm ($Y_1$)}  
		\end{table} 
	\end{landscape}
	
	\begin{landscape}
		\begin{table}[htp] 
			\centering 
			\begin{adjustbox}{height=0.55\textheight,width=1.6\textwidth}
				\begin{tabular}{c|ccc|ccc|ccc}
					\hline
					Model Name & 2h1 & 2h1\_ar & 2h1\_ar\_2sls & 2h1\_diss & 2h1\_diss\_ar & 2h1\_diss\_ar\_2sls & 2h1\_diss\_int & 2h1\_diss\_int\_ar & 2h1\_diss\_int\_ar\_2sls \\
					\hline
					Sample Size (N) & 44 & 44 & 25 & 44 & 44 & 25 & 44 & 44 & 25 \\
					\hline
					Endogenous &  &  & $\hat{Y_1}$ &  &  & $\hat{Y_1}$ &  &  & $\hat{Y_1}$ \\
					\hline
					C & 0.1156 & -0.0769 & -0.1483 & 0.2847 & 0.0915 & -0.1071 & 0.0189 & -0.1784 & -0.0572 \\
					& 0.1221 & 0.1285 & 0.2061 & 0.1275 & 0.1217 & 0.2416 & 0.1711 & 0.1519 & 0.2317 \\
					& 0.3496 & 0.5528 & 0.4807 & 0.0315** & 0.4571 & 0.6629 & 0.9128 & 0.2479 & 0.8078 \\
					\hline
					$Y_1$ &  & 0.6756 & 1.3902 &  & 0.7376 & 1.3499 &  & 0.7413 & 1.4680 \\
					&  & 0.2255 & 0.5230 &  & 0.1977 & 0.4384 &  & 0.1833 & 0.5081 \\
					&  & 0.0048*** & 0.0155** &  & 0.0006*** & 0.0065*** &  & 0.0003*** & 0.0102** \\
					\hline
					GENDER & 0.1431 & 0.0616 & -0.1182 & 0.1219 & 0.0308 & -0.1122 & 0.1377 & 0.0463 & -0.1345 \\
					& 0.0743 & 0.0730 & 0.1248 & 0.0689 & 0.0643 & 0.1133 & 0.0660 & 0.0599 & 0.1254 \\
					& 0.0615* & 0.4043 & 0.3557 & 0.0847* & 0.6349 & 0.3351 & 0.0439** & 0.4448 & 0.2983 \\
					Y$_2$GROUP & 0.1039 & 0.0712 & -0.0761 & 0.0789 & 0.0406 & -0.0726 & 0.0699 & 0.0313 & -0.0655 \\
					& 0.0687 & 0.0636 & 0.0767 & 0.0639 & 0.0562 & 0.0786 & 0.0610 & 0.0522 & 0.0817 \\
					& 0.1384 & 0.2698 & 0.3335 & 0.2245 & 0.4739 & 0.3677 & 0.2594 & 0.5520 & 0.4339 \\
					\hline
					DISS &  &  &  & -0.0597 & -0.0657 & -0.0140 & 0.0467 & 0.0419 & -0.0613 \\
					&  &  &  & 0.0212 & 0.0184 & 0.0455 & 0.0523 & 0.0440 & 0.0774 \\
					&  &  &  & 0.0076*** & 0.001*** & 0.7621 & 0.3770 & 0.3466 & 0.4395 \\
					PROXI$_{25}$ & -0.0945 & -0.1079 & 0.0654 & -0.0830 & -0.0965 & 0.0621 & -0.0283 & -0.0412 & 0.0850 \\
					& 0.0937 & 0.0855 & 0.1096 & 0.0865 & 0.0748 & 0.1111 & 0.0860 & 0.0724 & 0.1203 \\
					& 0.3196 & 0.2146 & 0.5577 & 0.3433 & 0.2050 & 0.5833 & 0.7444 & 0.5727 & 0.4895 \\
					PROXI$_{75}$ & 0.1282 & 0.1248 & 0.2723 & 0.1603 & 0.1598 & 0.2618 & 0.4878 & 0.4910 & 0.0819 \\
					& 0.0839 & 0.0765 & 0.0934 & 0.0782 & 0.0675 & 0.1007 & 0.1660 & 0.1396 & 0.2838 \\
					& 0.1347 & 0.1110 & 0.0089*** & 0.0473** & 0.0233** & 0.0181** & 0.0057*** & 0.0012*** & 0.7765 \\
					INT(DISS, PROXI$_{75}$) &  &  &  &  &  &  & -0.1322 & -0.1338 & 0.0901 \\
					&  &  &  &  &  &  & 0.0599 & 0.0504 & 0.1353 \\
					&  &  &  &  &  &  & 0.0336** & 0.0117** & 0.5145 \\
					\hline
					R\^{}2 & 0.2078 & 0.3591 & 0.5628 & 0.3447 & 0.5238 & 0.5741 & 0.4209 & 0.6018 & 0.5649 \\
					Adjusted R\^{}2 & 0.1265 & 0.2748 & 0.4477 & 0.2585 & 0.4466 & 0.4321 & 0.3270 & 0.5244 & 0.3857 \\
					Akaike info criterion & 0.0667 & -0.0999 & Not calculated & -0.0777 & -0.3514 & Not calculated & -0.1559 & -0.4848 & Not calculated \\
					Schwarz criterion & 0.2694 & 0.1434 & Not calculated & 0.1656 & -0.0676 & Not calculated & 0.1280 & -0.1604 & Not calculated \\
					Log likelihood & 3.5327 & 8.1983 & Not calculated & 7.7084 & 14.7316 & Not calculated & 10.4292 & 18.6666 & Not calculated \\
					\hline
					F-statistic & 2.5568 & 4.2592 & 3.8163 & 3.9980 & 6.7832 & 4.3463 & 4.4829 & 7.7722 & 3.7640 \\
					Prob(F-statistic) & 0.0538 & 0.0036 & 0.0146 & 0.0052 & 0.0001 & 0.0070 & 0.0017 & 0.0000 & 0.0121 \\
					J-statistic & NA & NA & 1.9436 & NA & NA & 1.8345 & NA & NA & 3.6342 \\
					Prob(J-statistic) & NA & NA & 0.5842 & NA & NA & 0.7662 & NA & NA & 0.6032 \\
					\hline
				\end{tabular}
			\end{adjustbox}
			\caption{Regression analysis for H1 for the second midterm ($Y_2$)}  
		\end{table} 
	\end{landscape}

	
	\subsection{H2 - General model for predicting performance}	
	H2 posited a linear relationship among attendance, homework “orderliness,” quiz performance, and perceived homework hardness: higher absence and lower homework organization would predict lower quiz scores and greater reported homework difficulty. The regression models reported in Table 4 do not support this hypothesis — coefficients were statistically insignificant, so H2 is not supported by our sample.\\
	
	These null results do not close the question. Key measurement limitations likely reduced power and increased noise. Homework in our dataset was not rubric-scored or coded as “sufficient/insufficient” despite some formative feedback to on‑time submitters; incorporating objective measures of homework quality would improve inference. Additional covariates that could yield insight include a student question-rate (as a proxy for engagement), lecture attention metrics, and counts of interactions during exercises. To avoid bias these should be operationalized reproducibly (for example, with standardized grading rubrics, automated participation logs, or consistent data collection across instructors). Greater granularity and more reliable measures would also allow exploration of non‑linear effects and interactions that our current linear models could not detect.
	
	\begin{landscape}
		\begin{table}[htp] 
			\centering 
			\begin{adjustbox}{height=0.5\textheight,width=1.6\textwidth}
			\begin{tabular}{c|ccc|ccc|cc|ccc}
				\hline
				Model Name & 1h2 & 1h2\_ar & 1h2\_ar\_2sls & 1h2\_reduced & 1h2\_reduced\_ar & 1h2\_reduced\_ar\_2sls & 2h2 & 2h2\_ar & 2h2\_reduced & 2h2\_reduced\_ar & 2h2\_reduced\_ar\_2sls \\
				\hline
				Depended & $Y_1$ & $Y_1$ & $Y_1$ & $Y_1$ & $Y_1$ & $Y_1$ & $Y_2$ & $Y_2$ & $Y_2$ & $Y_2$ & $Y_2$ \\
				\hline
				Sample Size (N) & 39 & 24 & 24 & 46 & 25 & 25 & 21 & 21 & 29 & 29 & 19 \\
				\hline
				Endogenous &  &  & $\hat{firstQIZZ}$ &  &  & $\hat{firstQIZZ}$ &  &  &  &  & $\hat{Y_1}$ \\
				\hline
				C & 0.5964 & 0.7843 & 0.8153 & 0.5634 & 0.8125 & 0.8109 & 0.0024 & -0.2451 & 0.0936 & -0.0371 & -0.1348 \\
				& 0.0999 & 0.1893 & 0.2285 & 0.0679 & 0.1497 & 0.1610 & 0.2633 & 0.1908 & 0.1823 & 0.1781 & 0.2240 \\
				& 0*** & 0.0009 & 0.0028*** & 0*** & 0*** & 0.0001*** & 0.9929 & 0.2232 & 0.6121 & 0.8369 & 0.5576 \\
				\hline
				firstQIZZ / $Y_1$ &  & -0.0079 & -0.5147 &  & 0.0248 & -0.2861 &  & 1.1376 &  & 0.6905 & 1.0741 \\
				&  & 0.2007 & 0.6003 &  & 0.1799 & 0.5424 &  & 0.2870 &  & 0.3058 & 0.6632 \\
				&  & 0.9690 & 0.4047 &  & 0.8919 & 0.6039 &  & 0.0019*** &  & 0.0337** & 0.1293 \\
				\hline
				GENDER & 0.1630 & 0.1313 & 0.1764 & 0.1425 & 0.0951 & 0.1250 & 0.2253 & -0.0411 & 0.2753 & 0.1613 & -0.0160 \\
				& 0.0526 & 0.0670 & 0.0938 & 0.0408 & 0.0587 & 0.0798 & 0.1091 & 0.1005 & 0.1045 & 0.1090 & 0.1681 \\
				& 0.0041*** & 0.069* & 0.0796* & 0.0012*** & 0.1219 & 0.1336 & 0.0595* & 0.6900 & 0.0145** & 0.1525 & 0.9256 \\
				Y$_{1or2}$GROUP & -0.1242 & -0.1596 & -0.1314 & -0.1124 & -0.1361 & -0.1206 & 0.0921 & 0.0246 & 0.0949 & 0.0325 & 0.1123 \\
				& 0.0499 & 0.0656 & 0.0841 & 0.0402 & 0.0553 & 0.0646 & 0.0997 & 0.0704 & 0.0964 & 0.0933 & 0.1089 \\
				& 0.0185** & 0.0279** & 0.1389 & 0.0079*** & 0.0236** & 0.0775* & 0.3723 & 0.7332 & 0.3347 & 0.7308 & 0.3214 \\
				\hline
				averATND\_H & -0.1680 & -0.2530 & -0.2196 & -0.1310 & -0.2898 & -0.2613 & -0.0335 & 0.0724 & 0.0880 & 0.0656 & 0.3883 \\
				& 0.0623 & 0.1467 & 0.1788 & 0.0497 & 0.1286 & 0.1459 & 0.1574 & 0.1110 & 0.1434 & 0.1329 & 0.2402 \\
				& 0.0112** & 0.1051 & 0.2383 & 0.0118** & 0.0362** & 0.0892* & 0.8350 & 0.5268 & 0.5451 & 0.6264 & 0.1300 \\
				lateHW & -0.0760 & 0.0135 & -0.0841 &  &  &  & -0.0096 & -0.0949 &  &  &  \\
				& 0.1191 & 0.1522 & 0.2104 &  &  &  & 0.1313 & 0.0925 &  &  &  \\
				& 0.5282 & 0.9307 & 0.6949 &  &  &  & 0.9430 & 0.3250 &  &  &  \\
				paperHW & -0.0197 & -0.0757 & -0.1035 &  &  &  & 0.2060 & 0.2213 &  &  &  \\
				& 0.0701 & 0.0924 & 0.1143 &  &  &  & 0.1475 & 0.1011 &  &  &  \\
				& 0.7810 & 0.4255 & 0.3798 &  &  &  & 0.1859 & 0.0491** &  &  &  \\
				averQIZZ\_H & 0.0509 & 0.0913 & 0.1185 &  &  &  & 0.0884 & 0.1347 &  &  &  \\
				& 0.0559 & 0.0636 & 0.0814 &  &  &  & 0.1144 & 0.0792 &  &  &  \\
				& 0.3700 & 0.1713 & 0.1661 &  &  &  & 0.4535 & 0.1149 &  &  &  \\
				hardHW\_H & 0.0833 & 0.0623 & 0.0854 & 0.0870 & 0.0514 & 0.0703 & -0.0300 & -0.1229 & 0.0842 & 0.0804 & -0.0954 \\
				& 0.0567 & 0.0733 & 0.0911 & 0.0460 & 0.0577 & 0.0693 & 0.1082 & 0.0777 & 0.1002 & 0.0927 & 0.1330 \\
				& 0.1519 & 0.4089 & 0.3633 & 0.0659* & 0.3841 & 0.3230 & 0.7857 & 0.1398 & 0.4092 & 0.3944 & 0.4859 \\
				\hline
				R\^{}2 & 0.4132 & 0.5900 & 0.4157 & 0.3924 & 0.5575 & 0.4880 & 0.5335 & 0.7980 & 0.2437 & 0.3810 & 0.5756 \\
				Adjusted R\^{}2 & 0.2807 & 0.3713 & 0.1041 & 0.3331 & 0.4411 & 0.3532 & 0.2823 & 0.6633 & 0.1177 & 0.2464 & 0.4123 \\
				Akaike info criterion & -0.9083 & -1.0721 & Not calculated & -1.0878 & -1.2417 & Not calculated & -0.0444 & -0.7860 & 0.1346 & 0.0033 & Not calculated \\
				Schwarz criterion & -0.5670 & -0.6303 & Not calculated & -0.8890 & -0.9492 & Not calculated & 0.3536 & -0.3384 & 0.3703 & 0.2862 & Not calculated \\
				Log likelihood & 25.7115 & 21.8654 & Not calculated & 30.0193 & 21.5210 & Not calculated & 8.4658 & 17.2531 & 3.0485 & 5.9520 & Not calculated \\
				\hline
				F-statistic & 3.1181 & 2.6982 & 1.9851 & 6.6183 & 4.7878 & 4.1899 & 2.1237 & 5.9247 & 1.9338 & 2.8311 & 2.6873 \\
				Prob(F-statistic) & 0.013064** & 0.046451** & 0.1201 & 0.000332*** & 0.005331*** & 0.0098*** & 0.1143 & 0.003237*** & 0.1374 & 0.039065** & 0.070116* \\
				J-statistic & NA & NA & 1.4442 & NA & NA & 2.0720 & NA & NA & NA & NA & 1.9637 \\
				Prob(J-statistic) & NA & NA & 0.6952 & NA & NA & 0.5576 & NA & NA & NA & NA & 0.3746 \\
				\hline
			\end{tabular}
	
			
			\end{adjustbox}
			\caption{Regression analysis for H2}  
			\end{table} 
	\end{landscape}
	
\section{Discussion}

This study examined how seating proximity, honesty, and classroom behaviours relate to student performance. Below are outlined the limitations, summarize key findings, and suggest directions for future research.

\subsection{Limitations}

Analyses rely on 49 observations from a single course cohort, with fewer complete cases for some specifications; this concentrated sample limits statistical power and external validity, and a substantial share of the cohort were non‑auditing, test‑only students, further narrowing the effective sample for richer inference. Scaling the design would be time‑ and labour‑intensive: several measures required timely collection and occasional re‑collection (for example, the evidence‑list honesty measure, \texttt{honLST}), which is practicable at small scale but burdensome for larger implementations and demands greater operational rigor from instructors and staff.\\

Key variables lack objective, fine‑grained measurement: homework quality was not rubric‑scored or coded for sufficiency/correctness, and the cheating proxy and honesty indicators rely partly on observable proxies rather than direct, objective measures. One individual recorded pre‑test behaviour, observed in‑test behaviour, and later participated in grading after an operational change; despite steps to minimise bias, this overlap may have produced observer effects and departures from the intended protocol.\\

Although the dishonesty index captured many overt actions, covert cheating and behaviour changes induced by awareness of observation remain possible; unobtrusive recording would be preferable but raises ethical and logistical constraints. Finally, prior mathematics achievement (for example, whether a student had passed prerequisite math and the grade obtained) was unavailable in primary specifications and is plausibly related to quiz performance; its exclusion may bias estimated relationships.

\subsection{Concluding summary}

Proximity to higher‑achieving peers (PROXI\textsubscript{75}) is robustly associated with higher end‑of‑course scores across OLS, autoregressive, and 2SLS specifications; this relationship persists when controlling for measured dishonesty and in the second midterm when collusion opportunities were reduced. Proximity to lower‑achieving peers (PROXI\textsubscript{25}) shows no consistent association. The dishonesty index (DISS) is a strong negative correlate of performance, on the order of approximately 7.6–11.2 percentage points per index point; by contrast, \texttt{atitMATH} does not predict outcomes in these models. The hypothesised strictly linear relationship among attendance, homework orderliness, quiz performance, and perceived homework difficulty (H2) is not supported by the current sample.

\subsection{Recommendations for future research}

Replicate the design with larger and more heterogeneous samples across multiple cohorts, courses and institutions to improve power and external validity. Adopt rubric‑based homework scoring, collect objective prior‑ability indicators (for example, prior math grades or standardized test scores), and capture automated participation or attention metrics where feasible to improve measurement.\\

Strengthen protocols to mitigate observer effects by separating data collection from grading, pre‑registering procedures, and—where ethical—using independent proctoring or automated monitoring to reduce awareness biases. Address dishonesty explicitly through robustness checks that exclude or instrument for flagged observations and by incorporating validated behavioural measures of cheating. Broaden modelling strategies to test non‑linearities and interactions (for example, between proximity and ability) and to employ panel or mixed‑effects models with longitudinal data to better isolate causal mechanisms.\\

Taken together, the findings point to substantive relationships—notably the positive association of proximity to higher‑achieving peers and the negative association of dishonesty—that merit further investigation; however, larger, better‑measured samples and tighter protocols are required before asserting definitive causal claims.
	
	\begin{thebibliography}{99}
		
		\bibitem{anderman2007}
		Anderman, E. M., \& Murdock, T. B. (2007). \textit{Psychology of academic cheating}. Elsevier. ISBN: 978-0-12-372541-7. \href{https://isbnsearch.org/isbn/9780123725417}{https://isbnsearch.org/isbn/9780123725417}
		
		\bibitem{balfanz2012}
		Balfanz, R., \& Byrnes, V. (2012). \textit{The importance of being in school: A report on absenteeism in the nation's public schools}. Johns Hopkins University Center for Social Organization of Schools.out of 49  \href{https://new.every1graduates.org/wp-content/uploads/2012/05/FINALChronicAbsenteeismReport_May16.pdf}{https://new.every1graduates.org/wp-content/uploads/2012/05/FINALChronicAbsenteeismReport\_May16.pdf}
		
		\bibitem{black1998}
		Black, P., \& Wiliam, D. (1998). Assessment and classroom learning. \textit{Assessment in Education: Principles, Policy \& Practice}, 5(1), 7--74. \href{https://doi.org/10.1080/0969595980050102}{https://doi.org/10.1080/0969595980050102}
		
		\bibitem{carrell2013}
		Carrell, S. E., Sacerdote, B., \& West, J. E. (2013). From natural variation to optimal policy? The importance of endogenous peer group formation. \textit{Econometrica}, 81(3), 855--882. \href{https://www.jstor.org/stable/23524165}{https://www.jstor.org/stable/23524165}
		
		\bibitem{cooper2006}
		Cooper, H., Robinson, J. C., \& Patall, E. A. (2006). Does homework improve academic achievement? A synthesis of research, 1987--2003. \textit{Review of Educational Research}, 76(1), 1--62. \href{https://doi.org/10.3102/00346543076001001}{https://doi.org/10.3102/00346543076001001}
		
		\bibitem{gottfried2010}
		Gottfried, M. A. (2010). Evaluating the relationship between student attendance and achievement in urban elementary and middle schools: An instrumental variables approach. \textit{American Educational Research Journal}, 47(2), 434--465. \href{https://doi.org/10.3102/0002831209350494}{https://doi.org/10.3102/0002831209350494}
		
		\bibitem{hattie2007}
		Hattie, J., \& Timperley, H. (2007). The power of feedback. \textit{Review of Educational Research}, 77(1), 81--112. \href{https://doi.org/10.3102/003465430298487}{https://doi.org/10.3102/003465430298487}
		
		\bibitem{konig2020}
		König, J., Jäger-Biela, D. J., \& Glutsch, N. (2020). Adapting to online teaching during COVID-19 school closure: Teacher education and teacher competence effects among German preservice teachers. \textit{European Journal of Teacher Education}, 43(4), 608--622. \href{https://doi.org/10.1080/02619768.2020.1809650}{https://doi.org/10.1080/02619768.2020.1809650}
		
		\bibitem{mccabe2012}
		McCabe, D. L., Butterfield, K. D., \& Treviño, L. K. (2012). \textit{Cheating in college: Why students do it and what educators can do about it}. Johns Hopkins University Press. \href{https://jhupbooks.press.jhu.edu/title/cheating-college}{https://jhupbooks.press.jhu.edu/title/cheating-college}
		
		\bibitem{pajares1999}
		Pajares, F., \& Graham, L. (1999). Self-efficacy, motivation constructs, and mathematics performance of entering middle school students. \textit{Contemporary Educational Psychology}, 24(2), 124--139. \href{https://doi.org/10.1006/ceps.1998.0991}{https://doi.org/10.1006/ceps.1998.0991}
		
		\bibitem{sacerdote2011}
		Sacerdote, B. (2011). Peer effects in education: How might they work, how big are they and how much do we know thus far? In E. A. Hanushek, S. Machin, \& L. Woessmann (Eds.), \textit{Handbook of the Economics of Education} (Vol. 3, pp. 249--277). Elsevier. \href{https://doi.org/10.1016/B978-0-444-53429-3.00004-1}{https://doi.org/10.1016/B978-0-444-53429-3.00004-1}
		
		\bibitem{trautwein2009}
		Trautwein, U., Lüdtke, O., Köller, O., \& Baumert, J. (2009). Self-esteem, academic self-concept, and achievement: How the learning environment moderates their relations. \textit{Contemporary Educational Psychology}, 34(2), 163--175. \href{https://doi.org/10.1016/j.cedpsych.2009.01.002}{https://doi.org/10.1016/j.cedpsych.2009.01.002}
		
		\bibitem{zimmerman2000}
		Zimmerman, B. J. (2000). Self-efficacy: An essential motive to learn. \textit{Contemporary Educational Psychology}, 25(1), 82--91. \href{https://doi.org/10.1006/ceps.1999.1016}{https://doi.org/10.1006/ceps.1999.1016}
		
	\end{thebibliography}
	
	
	
\end{document}
